\documentclass[12pt]{article}
\usepackage{amsmath, amssymb, geometry}
\geometry{a4paper, margin=1in}
\usepackage{parskip}

\title{Homework 4: Homogeneous Linear Equations: The General Solution}
\author{Tashfeen Omran}
\date{September 2026}

\begin{document}

\maketitle
\tableofcontents
\newpage

\section{Part 1: Comprehensive Introduction, Context, and Prerequisites}

\subsection{Core Concepts}
In this chapter, we begin our study of linear second-order constant-coefficient differential equations. The general form of these equations is:
\[ a y'' + b y' + c y = f(t) \]
where $a$, $b$, and $c$ are real constants and $a \neq 0$. When the forcing function $f(t)$ is identically zero, the equation is called \textbf{homogeneous}:
\[ a y'' + b y' + c y = 0 \]
The core concept here is the \textbf{Principle of Superposition}. Because the equation is linear, if we can find two distinct (linearly independent) solutions, $y_1(t)$ and $y_2(t)$, any linear combination of these solutions will also be a solution. Thus, the general solution is formed as:
\[ y(t) = c_1 y_1(t) + c_2 y_2(t) \]
where $c_1$ and $c_2$ are arbitrary constants determined by initial conditions.

\subsection{Intuition and Derivation}
To solve $a y'' + b y' + c y = 0$, we need a function whose derivatives are constant multiples of the function itself, so that they can cancel out when added together. The exponential function $y(t) = e^{rt}$ perfectly fits this description. 
Substituting $y = e^{rt}$, $y' = r e^{rt}$, and $y'' = r^2 e^{rt}$ into the differential equation yields:
\[ a(r^2 e^{rt}) + b(r e^{rt}) + c(e^{rt}) = 0 \]
Factoring out $e^{rt}$ (which is never zero), we are left with the \textbf{auxiliary equation} (or characteristic equation):
\[ a r^2 + b r + c = 0 \]
This reduces the calculus problem of solving a differential equation into an algebra problem of finding the roots of a quadratic polynomial.

\subsection{Historical Context and Motivation}
The study of second-order linear differential equations blossomed in the 18th century, driven by the need to understand classical mechanics. Mathematicians and physicists like Leonhard Euler and Daniel Bernoulli were studying the vibrations of strings, the motion of pendulums, and mass-spring systems. These physical systems naturally exhibit acceleration (second derivative) proportional to position and velocity, leading directly to equations of the form $m y'' + b y' + k y = 0$. The algebraic approach using the characteristic equation was largely formalized by Euler, providing a systematic way to analyze mechanical vibrations.

\subsection{Key Formulas}
\begin{itemize}
    \item \textbf{Auxiliary Equation:} $a r^2 + b r + c = 0$
    \item \textbf{Distinct Real Roots ($r_1 \neq r_2$):} If the discriminant $b^2 - 4ac > 0$, the general solution is:
    \[ y(t) = c_1 e^{r_1 t} + c_2 e^{r_2 t} \]
    \item \textbf{Repeated Real Root ($r_1 = r_2 = r$):} If the discriminant $b^2 - 4ac = 0$, we need a second linearly independent solution. We multiply by $t$ to get:
    \[ y(t) = c_1 e^{rt} + c_2 t e^{rt} \]
    \item \textbf{The Wronskian:} A determinant used to test for linear independence of two functions $y_1$ and $y_2$:
    \[ W[y_1, y_2](t) = \det \begin{pmatrix} y_1(t) & y_2(t) \\ y_1'(t) & y_2'(t) \end{pmatrix} = y_1(t) y_2'(t) - y_1'(t) y_2(t) \]
    If $W \neq 0$, the functions are linearly independent.
\end{itemize}

\subsection{Prerequisites}
To succeed in this topic, you must be proficient in:
\begin{enumerate}
    \item \textbf{Algebra:} Factoring quadratic and cubic polynomials, using the quadratic formula, and solving systems of linear equations.
    \item \textbf{Calculus:} Basic differentiation rules, especially the chain rule and the product rule (crucial for differentiating $t e^{rt}$).
\end{enumerate}

\newpage
\section{Part 2: Detailed Homework Solutions}
\textit{Note: A strategic selection of problems has been made to cover every distinct problem archetype (Distinct Roots, Repeated Roots, Initial Value Problems, Boundary Value Problems, Linear Independence, and Higher-Order Equations) without redundant repetition, maximizing learning efficiency.}

\subsection{Archetype 1: General Solution, Distinct Real Roots}
\textbf{Problem 1:} Find a general solution to the given differential equation: $2y'' + 7y' - 4y = 0$.

\textbf{Solution:}
Step 1: Write the auxiliary equation.
\[ 2r^2 + 7r - 4 = 0 \]
Step 2: Find the roots of the auxiliary equation by factoring.
We look for two numbers that multiply to $(2)(-4) = -8$ and add to $7$. These numbers are $8$ and $-1$.
\[ 2r^2 + 8r - r - 4 = 0 \]
\[ 2r(r + 4) - 1(r + 4) = 0 \]
\[ (2r - 1)(r + 4) = 0 \]
The roots are $r_1 = \frac{1}{2}$ and $r_2 = -4$.
Step 3: Write the general solution. Since the roots are real and distinct, the solution is:
\[ y(t) = c_1 e^{t/2} + c_2 e^{-4t} \]

\subsection{Archetype 2: General Solution, Repeated Real Roots}
\textbf{Problem 11:} Find a general solution to the given differential equation: $4w'' + 20w' + 25w = 0$.

\textbf{Solution:}
Step 1: Write the auxiliary equation.
\[ 4r^2 + 20r + 25 = 0 \]
Step 2: Find the roots. Notice that this is a perfect square trinomial.
\[ (2r + 5)^2 = 0 \]
The root is $r = -\frac{5}{2}$, which is a repeated root (multiplicity 2).
Step 3: Write the general solution. For a repeated root $r$, the linearly independent solutions are $e^{rt}$ and $t e^{rt}$.
\[ w(t) = c_1 e^{-5t/2} + c_2 t e^{-5t/2} \]

\subsection{Archetype 3: Initial Value Problem (Distinct Roots)}
\textbf{Problem 13:} Solve the initial value problem: $y'' + 2y' - 8y = 0; \quad y(0) = 3, \quad y'(0) = -12$.

\textbf{Solution:}
Step 1: Find the general solution. The auxiliary equation is:
\[ r^2 + 2r - 8 = 0 \]
\[ (r + 4)(r - 2) = 0 \implies r_1 = -4, \quad r_2 = 2 \]
The general solution is:
\[ y(t) = c_1 e^{-4t} + c_2 e^{2t} \]
Step 2: Differentiate the general solution to use the second initial condition.
\[ y'(t) = -4c_1 e^{-4t} + 2c_2 e^{2t} \]
Step 3: Apply the initial conditions $y(0) = 3$ and $y'(0) = -12$.
\begin{align*}
    y(0) &= c_1 e^0 + c_2 e^0 = c_1 + c_2 = 3 \\
    y'(0) &= -4c_1 e^0 + 2c_2 e^0 = -4c_1 + 2c_2 = -12
\end{align*}
Step 4: Solve the system of linear equations.
From the first equation, $c_2 = 3 - c_1$. Substitute this into the second equation:
\begin{align*}
    -4c_1 + 2(3 - c_1) &= -12 \\
    -4c_1 + 6 - 2c_1 &= -12 \\
    -6c_1 &= -18 \\
    c_1 &= 3
\end{align*}
Now find $c_2$:
\[ c_2 = 3 - 3 = 0 \]
Step 5: Write the final specific solution.
\[ y(t) = 3 e^{-4t} + 0 e^{2t} \implies y(t) = 3 e^{-4t} \]

\subsection{Archetype 4: Initial Value Problem (Repeated Roots)}
\textbf{Problem 19:} Solve the initial value problem: $y'' + 2y' + y = 0; \quad y(0) = 1, \quad y'(0) = -3$.

\textbf{Solution:}
Step 1: Find the general solution. The auxiliary equation is:
\[ r^2 + 2r + 1 = 0 \implies (r + 1)^2 = 0 \implies r = -1 \text{ (repeated)} \]
The general solution is:
\[ y(t) = c_1 e^{-t} + c_2 t e^{-t} \]
Step 2: Differentiate the general solution. \textbf{Caution:} Use the product rule for the second term!
\[ y'(t) = -c_1 e^{-t} + c_2 \left( 1 \cdot e^{-t} + t \cdot (-e^{-t}) \right) = -c_1 e^{-t} + c_2 e^{-t} - c_2 t e^{-t} \]
Step 3: Apply the initial conditions.
\begin{align*}
    y(0) &= c_1 e^0 + c_2(0)e^0 = c_1 = 1 \\
    y'(0) &= -c_1 e^0 + c_2 e^0 - c_2(0)e^0 = -c_1 + c_2 = -3
\end{align*}
Step 4: Solve for the constants.
We already have $c_1 = 1$. Substitute into the second equation:
\[ -1 + c_2 = -3 \implies c_2 = -2 \]
Step 5: Write the final specific solution.
\[ y(t) = e^{-t} - 2t e^{-t} \]

\subsection{Archetype 5: Boundary Value Problems vs. Initial Value Problems}
\textbf{Problem 26:} Given that every solution to $y'' + y = 0$ is of the form $y(t) = c_1 \cos t + c_2 \sin t$, show that for boundary value problems there is no existence-uniqueness theorem analogous to Theorem 1.
(a) Show there is a unique solution satisfying $y(0) = 2$ and $y(\pi/2) = 0$.
(b) Show there is no solution satisfying $y(0) = 2$ and $y(\pi) = 0$.
(c) Show there are infinitely many solutions satisfying $y(0) = 2$ and $y(\pi) = -2$.

\textbf{Solution:}
\textbf{Part (a):} Apply the boundary conditions to $y(t) = c_1 \cos t + c_2 \sin t$.
\begin{align*}
    y(0) &= c_1 \cos(0) + c_2 \sin(0) = c_1 = 2 \\
    y(\pi/2) &= c_1 \cos(\pi/2) + c_2 \sin(\pi/2) = c_2 = 0
\end{align*}
Thus, $c_1 = 2$ and $c_2 = 0$. The unique solution is $y(t) = 2 \cos t$.

\textbf{Part (b):} Apply the boundary conditions.
\begin{align*}
    y(0) &= c_1 \cos(0) + c_2 \sin(0) = c_1 = 2 \\
    y(\pi) &= c_1 \cos(\pi) + c_2 \sin(\pi) = -c_1 = 0
\end{align*}
This implies $c_1 = 2$ and $-c_1 = 0$, which is a contradiction ($2 \neq 0$). Therefore, \textbf{no solution exists}.

\textbf{Part (c):} Apply the boundary conditions.
\begin{align*}
    y(0) &= c_1 \cos(0) + c_2 \sin(0) = c_1 = 2 \\
    y(\pi) &= c_1 \cos(\pi) + c_2 \sin(\pi) = -c_1 = -2
\end{align*}
Both conditions yield $c_1 = 2$. The constant $c_2$ is completely unrestricted. Therefore, any function of the form $y(t) = 2 \cos t + c_2 \sin t$ is a solution. There are \textbf{infinitely many solutions}.

\subsection{Archetype 6: Linear Independence and the Wronskian}
\textbf{Problem 27:} Determine whether the functions $y_1(t) = \cos t \sin t$ and $y_2(t) = \sin 2t$ are linearly dependent or independent on the interval $(0, 1)$.

\textbf{Solution:}
Two functions are linearly dependent if one is a constant multiple of the other. Recall the double-angle trigonometric identity:
\[ \sin 2t = 2 \sin t \cos t \]
Therefore, we can write:
\[ y_2(t) = 2 y_1(t) \]
Since $y_2(t)$ is a constant multiple of $y_1(t)$, the functions are \textbf{linearly dependent}.

\textbf{Problem 34(a):} Show that $W[y_1, y_2](t) = y_1(t)y_2'(t) - y_1'(t)y_2(t)$ can be conveniently expressed as the $2 \times 2$ determinant.

\textbf{Solution:}
The determinant of a $2 \times 2$ matrix $\begin{pmatrix} a & b \\ c & d \end{pmatrix}$ is defined as $ad - bc$.
Setting up the Wronskian matrix with the functions in the first row and their derivatives in the second row:
\[ \det \begin{pmatrix} y_1(t) & y_2(t) \\ y_1'(t) & y_2'(t) \end{pmatrix} = y_1(t) \cdot y_2'(t) - y_2(t) \cdot y_1'(t) \]
This perfectly matches the definition of the Wronskian $W[y_1, y_2](t)$.

\subsection{Archetype 7: Higher-Order Linear Equations}
\textbf{Problem 43:} Solve the initial value problem: $y''' - y' = 0; \quad y(0) = 2, \quad y'(0) = 3, \quad y''(0) = -1$.

\textbf{Solution:}
Step 1: Write the auxiliary equation for this third-order differential equation.
\[ r^3 - r = 0 \]
Step 2: Factor to find the roots.
\[ r(r^2 - 1) = 0 \implies r(r - 1)(r + 1) = 0 \]
The roots are $r_1 = 0$, $r_2 = 1$, and $r_3 = -1$.
Step 3: Write the general solution.
\[ y(t) = c_1 e^{0t} + c_2 e^{1t} + c_3 e^{-1t} = c_1 + c_2 e^t + c_3 e^{-t} \]
Step 4: Find the first and second derivatives.
\begin{align*}
    y'(t) &= c_2 e^t - c_3 e^{-t} \\
    y''(t) &= c_2 e^t + c_3 e^{-t}
\end{align*}
Step 5: Apply the initial conditions to form a $3 \times 3$ system.
\begin{align*}
    y(0) &= c_1 + c_2 + c_3 = 2 \\
    y'(0) &= c_2 - c_3 = 3 \\
    y''(0) &= c_2 + c_3 = -1
\end{align*}
Step 6: Solve the system. Add the second and third equations:
\[ (c_2 - c_3) + (c_2 + c_3) = 3 + (-1) \implies 2c_2 = 2 \implies c_2 = 1 \]
Substitute $c_2 = 1$ into the third equation:
\[ 1 + c_3 = -1 \implies c_3 = -2 \]
Substitute $c_2 = 1$ and $c_3 = -2$ into the first equation:
\[ c_1 + 1 - 2 = 2 \implies c_1 - 1 = 2 \implies c_1 = 3 \]
Step 7: Write the final solution.
\[ y(t) = 3 + e^t - 2e^{-t} \]

\newpage
\section{Part 3: In-Depth Analysis of Problems and Techniques}

\subsection{A) Problem Types and General Approach}
\begin{itemize}
    \item \textbf{General Solution Problems (Distinct/Repeated Roots):} The strategy is purely algebraic. Convert the ODE to its auxiliary polynomial, find the roots, and map them to the corresponding exponential forms. If a root $r$ repeats, multiply the second term by $t$ to ensure linear independence.
    \item \textbf{Initial Value Problems (IVPs):} These require finding the general solution first, then differentiating it as many times as necessary (one less than the order of the ODE). Finally, substitute $t=0$ (or $t=t_0$) to create a system of linear equations to solve for $c_1, c_2, \dots, c_n$.
    \item \textbf{Boundary Value Problems (BVPs):} Unlike IVPs where all conditions are at the same point $t_0$, BVPs have conditions at different points (e.g., $t=0$ and $t=\pi$). The strategy is the same as IVPs, but as seen in Problem 26, BVPs do not guarantee a unique solution; they can have zero, one, or infinitely many solutions.
    \item \textbf{Higher-Order Equations:} The strategy scales naturally. An $n$-th order ODE yields an $n$-th degree polynomial. You must find $n$ roots and solve an $n \times n$ system of equations for IVPs.
\end{itemize}

\subsection{B) Key Algebraic and Calculus Manipulations}
\begin{itemize}
    \item \textbf{Factoring Quadratics and Cubics:} Essential for finding roots. In Problem 43, factoring out the greatest common factor $r$ from $r^3 - r = 0$ was the crucial first step to avoid losing the root $r=0$.
    \item \textbf{The Product Rule for Repeated Roots:} In Problem 19, differentiating $c_2 t e^{-t}$ requires the product rule: $\frac{d}{dt}[t e^{-t}] = 1 \cdot e^{-t} + t \cdot (-e^{-t})$. Forgetting this is the most common calculus error in this chapter.
    \item \textbf{Solving Linear Systems via Elimination:} In Problem 13 and 43, adding equations together to eliminate variables (e.g., adding $c_2 - c_3 = 3$ and $c_2 + c_3 = -1$ to eliminate $c_3$) is much faster and less error-prone than substitution.
\end{itemize}

\newpage
\section{Part 4: "Cheatsheet" and Tips for Success}

\subsection{Summary of Important Formulas}
\begin{itemize}
    \item \textbf{ODE:} $a y'' + b y' + c y = 0$
    \item \textbf{Auxiliary:} $a r^2 + b r + c = 0$
    \item \textbf{Roots $r_1 \neq r_2$:} $y(t) = c_1 e^{r_1 t} + c_2 e^{r_2 t}$
    \item \textbf{Roots $r_1 = r_2 = r$:} $y(t) = c_1 e^{rt} + c_2 t e^{rt}$
    \item \textbf{Wronskian:} $W = y_1 y_2' - y_1' y_2$
\end{itemize}

\subsection{Tricks and Shortcuts}
\begin{itemize}
    \item \textbf{The $e^0$ Shortcut:} When solving IVPs at $t=0$, remember that $e^0 = 1$. You can often write down the linear system immediately without writing out the full $e^{rt}$ terms. For $y(t) = c_1 e^{r_1 t} + c_2 e^{r_2 t}$, the system is always:
    \begin{align*}
        c_1 + c_2 &= y(0) \\
        r_1 c_1 + r_2 c_2 &= y'(0)
    \end{align*}
    \item \textbf{Missing $y'$ Term:} If the ODE is $y'' - k^2 y = 0$, the roots are simply $r = \pm k$.
\end{itemize}

\subsection{Common Pitfalls and Traps}
\begin{itemize}
    \item \textbf{Forgetting the $t$:} Writing $y(t) = c_1 e^{rt} + c_2 e^{rt}$ for a repeated root. This collapses into a single constant $(c_1+c_2)e^{rt}$ and is not a general solution. You \textit{must} include the $t$.
    \item \textbf{Sign Errors in the Auxiliary Equation:} Translating $y'' - 4y = 0$ to $r^2 + 4 = 0$ instead of $r^2 - 4 = 0$. Always double-check your signs.
    \item \textbf{Confusing IVPs and BVPs:} Assuming a BVP will always have a unique solution. Always check for contradictions (like $2 = 0$) which indicate no solution.
\end{itemize}

\newpage
\section{Part 5: Conceptual Synthesis and The "Big Picture"}

\subsection{A) Thematic Connections}
The central theme of this topic is \textbf{Linearity and Superposition}. This is a profound concept that appears throughout mathematics, particularly in Linear Algebra. The set of all solutions to a homogeneous linear differential equation forms a \textit{vector space}. The functions $e^{r_1 t}$ and $e^{r_2 t}$ act as \textit{basis vectors}. Just as any point in 2D space can be reached by a linear combination of the $\hat{i}$ and $\hat{j}$ vectors ($x\hat{i} + y\hat{j}$), any solution to the differential equation can be reached by a linear combination of the basis functions ($c_1 y_1 + c_2 y_2$). The Wronskian is essentially a test to ensure your basis vectors are not pointing in the same direction (i.e., they are linearly independent).

\subsection{B) Forward and Backward Links}
\textbf{Backward Links:} This topic builds directly on first-order linear equations. A first-order equation $y' - ry = 0$ has the solution $y = c e^{rt}$. The second-order theory is a logical extension: we guess the same exponential form, which leads to a polynomial of degree 2 instead of degree 1.

\textbf{Forward Links:} This is the absolute bedrock for the rest of differential equations. 
1. \textbf{Complex Roots:} In the very next section, you will encounter $b^2 - 4ac < 0$. The algebraic roots will be complex ($r = \alpha \pm i\beta$), and you will use Euler's Formula ($e^{i\theta} = \cos\theta + i\sin\theta$) to translate these into real-valued oscillating solutions involving sines and cosines.
2. \textbf{Non-Homogeneous Equations:} When the equation equals a forcing function $f(t)$ instead of $0$, you will first find the homogeneous solution $y_h$ (exactly what we did here), and then add a particular solution $y_p$ to get the full solution $y = y_h + y_p$.

\newpage
\section{Part 6: Real-World Application and Modeling}

\subsection{A) Concrete Scenarios in Finance and Economics}
\begin{enumerate}
    \item \textbf{Macroeconomic Business Cycles (The Samuelson Oscillator):} In economics, the interaction between the multiplier effect (consumer spending) and the accelerator effect (capital investment) can cause cyclical fluctuations in a nation's Gross Domestic Product (GDP). When modeled in continuous time, this dynamic system reduces to a second-order linear homogeneous differential equation. The roots of the auxiliary equation determine whether the economy will experience runaway inflation, steady oscillation, or exponential decay back to equilibrium.
    \item \textbf{Quantitative Finance (Transforming the Black-Scholes PDE):} The famous Black-Scholes equation for pricing European options is a parabolic partial differential equation. Through a series of strategic variable substitutions (changing variables to dimensionless moneyness and time to maturity), financial mathematicians transform the complex Black-Scholes PDE into the standard Heat Equation, which is a linear constant-coefficient equation. Solving this relies heavily on the foundational principles of linear combinations of exponential functions.
\end{enumerate}

\subsection{B) Model Problem Setup: The Continuous-Time Samuelson Model}
Let $Y(t)$ represent the national income (GDP) at time $t$. In a simplified continuous-time multiplier-accelerator model, the rate of change of income is driven by past income levels and the rate of change of investment. 
Through economic derivation, the system can be modeled by the second-order equation:
\[ Y''(t) + \alpha Y'(t) + \beta Y(t) = 0 \]
where:
\begin{itemize}
    \item $\alpha$ is a constant related to the marginal propensity to save and the speed of market adjustment.
    \item $\beta$ is a constant related to the capital-output ratio (the accelerator).
\end{itemize}
To determine the long-term behavior of the economy, an economist would need to solve the auxiliary equation:
\[ r^2 + \alpha r + \beta = 0 \]
If the roots are real and negative, the economy stabilizes. If the roots are complex (which we will study next), the economy experiences boom-and-bust business cycles.

\newpage
\section{Part 7: Common Variations and Untested Concepts}

\subsection{Missing Concept 1: Complex Roots (Preview)}
While the homework focused on real roots, the most common variation of $a r^2 + b r + c = 0$ is when $b^2 - 4ac < 0$, resulting in complex roots $r = \alpha \pm i\beta$. 
\textbf{Explanation:} Using Euler's formula, the complex exponentials are converted into real-valued trigonometric functions. The general solution becomes:
\[ y(t) = e^{\alpha t} (c_1 \cos(\beta t) + c_2 \sin(\beta t)) \]
\textbf{Example:} Solve $y'' + 9y = 0$.
The auxiliary equation is $r^2 + 9 = 0 \implies r^2 = -9 \implies r = \pm 3i$.
Here, $\alpha = 0$ and $\beta = 3$. The solution is:
\[ y(t) = e^{0t} (c_1 \cos(3t) + c_2 \sin(3t)) = c_1 \cos(3t) + c_2 \sin(3t) \]

\subsection{Missing Concept 2: Reduction of Order}
The homework tells you to multiply by $t$ for repeated roots, but it doesn't test \textit{why} this works. The technique used to prove this is called Reduction of Order.
\textbf{Explanation:} If you know one solution $y_1(t)$, you can find a second linearly independent solution by guessing $y_2(t) = v(t)y_1(t)$, where $v(t)$ is an unknown function. Substituting this into the ODE reduces the second-order equation to a first-order equation for $v'(t)$.
\textbf{Example:} For $y'' - 2y' + y = 0$, the root is $r=1$, so $y_1 = e^t$.
Guess $y_2 = v(t)e^t$.
$y_2' = v'e^t + ve^t$
$y_2'' = v''e^t + 2v'e^t + ve^t$
Substitute into ODE:
$(v''e^t + 2v'e^t + ve^t) - 2(v'e^t + ve^t) + (ve^t) = 0$
$v''e^t = 0 \implies v'' = 0 \implies v(t) = c_1 t + c_2$.
Choosing $c_1=1, c_2=0$ gives $v(t) = t$, proving that $y_2(t) = t e^t$ is the second solution!

\newpage
\section{Part 8: Advanced Diagnostic Testing: "Find the Flaw"}
\textit{Instructions: You are the grader. Each of the following five solutions contains a single, subtle mathematical or conceptual error. Identify the error, explain why it is wrong in one sentence, and provide the correct final solution.}

\subsection{Problem 1}
\textbf{Question:} Solve the differential equation $y'' - 5y' - 6y = 0$.
\textbf{Flawed Solution:}
The auxiliary equation is $r^2 - 5r - 6 = 0$.
Factoring yields $(r - 2)(r - 3) = 0$.
The roots are $r = 2$ and $r = 3$.
The general solution is $y(t) = c_1 e^{2t} + c_2 e^{3t}$.

\subsection{Problem 2}
\textbf{Question:} Solve the initial value problem $y'' - 6y' + 9y = 0; \quad y(0)=2, \quad y'(0)=5$.
\textbf{Flawed Solution:}
The auxiliary equation is $r^2 - 6r + 9 = 0 \implies (r-3)^2 = 0$.
The repeated root is $r=3$.
The general solution is $y(t) = c_1 e^{3t} + c_2 t e^{3t}$.
Applying $y(0)=2$: $c_1 e^0 + c_2(0)e^0 = 2 \implies c_1 = 2$.
Differentiating: $y'(t) = 3c_1 e^{3t} + 3c_2 t e^{3t}$.
Applying $y'(0)=5$: $3(2)e^0 + 3c_2(0)e^0 = 5 \implies 6 = 5$.
This is a contradiction, so no solution exists.

\subsection{Problem 3}
\textbf{Question:} Find the general solution to $y''' - 4y' = 0$.
\textbf{Flawed Solution:}
The auxiliary equation is $r^3 - 4r = 0$.
Divide both sides by $r$: $r^2 - 4 = 0$.
$r^2 = 4 \implies r = 2, -2$.
The general solution is $y(t) = c_1 e^{2t} + c_2 e^{-2t}$.

\subsection{Problem 4}
\textbf{Question:} Calculate the Wronskian of $y_1(t) = e^{2t}$ and $y_2(t) = e^{-2t}$ at $t=0$.
\textbf{Flawed Solution:}
$y_1'(t) = 2e^{2t}$ and $y_2'(t) = -2e^{-2t}$.
The Wronskian is $W = y_1 y_2' + y_1' y_2$.
$W = (e^{2t})(-2e^{-2t}) + (2e^{2t})(e^{-2t})$.
$W = -2e^0 + 2e^0 = -2 + 2 = 0$.
Since $W=0$, the functions are linearly dependent.

\subsection{Problem 5}
\textbf{Question:} Solve the boundary value problem $y'' - y = 0; \quad y(0)=1, \quad y(1)=e$.
\textbf{Flawed Solution:}
Auxiliary equation: $r^2 - 1 = 0 \implies r = \pm 1$.
General solution: $y(t) = c_1 e^t + c_2 e^{-t}$.
$y(0) = c_1 + c_2 = 1 \implies c_2 = 1 - c_1$.
$y(1) = c_1 e^1 + c_2 e^{-1} = e$.
Substitute $c_2$: $c_1 e + (1 - c_1)e^{-1} = e$.
$c_1 e + e^{-1} - c_1 e^{-1} = e$.
$c_1(e - e^{-1}) = e - e^{-1}$.
Divide by $(e - e^{-1})$: $c_1 = 1$.
Then $c_2 = 1 - 1 = 0$.
The solution is $y(t) = 1 \cdot e^t + 0 \cdot e^{-t} = e^t$.
Wait, let me check $y(1) = e^1 = e$. It works. The solution is $y(t) = e^t$.
\textit{(Hint for the grader: Look closely at the algebraic step where we divide by $(e - e^{-1})$. Is there a conceptual flaw in how BVPs are treated, or did the student get lucky? Actually, look at the roots again. Are they correct? Yes. Look at the algebra. Is it correct? Yes. Wait, is there a flaw? Read carefully: The flaw is in the assumption that dividing by $(e - e^{-1})$ is the ONLY way to solve it, or perhaps there is a missing solution? No, linear BVPs can have unique solutions. Let's provide a real flaw.)}

\textbf{Revised Flawed Solution for Problem 5:}
Auxiliary equation: $r^2 - 1 = 0 \implies r = 1$. (Forgot the $\pm$).
General solution: $y(t) = c_1 e^t$.
$y(0) = c_1 e^0 = 1 \implies c_1 = 1$.
$y(1) = 1 \cdot e^1 = e$.
The solution is $y(t) = e^t$.

\end{document}